\section{Numerical accuracy of the headline kernel}
\label{sec:accuracy}

A kernel that is fast and wrong is not a result, and a tolerance chosen after the
fact is not a gate. This section states the error model for the top FP16 kernel,
what it predicts, and what the test suite measures against it.

\subsection{The model}

The kernel has two distinct sources of rounding and they do not have the same
size.

Inside one \texttt{mma.sync.aligned.m16n8k16} the sixteen products of a K slice
are summed by the hardware, in FP32, in a fixed reduction network. That step
contributes on the order of one FP32 rounding, not sixteen, and its ordering is a
property of the instruction rather than of my code~\cite{markidis2018}.

Across K the kernel accumulates those per instruction results sequentially in
FP32, one addition per \texttt{mma.sync}, so a contraction of length $k$ carries
$k/16$ dependent FP32 roundings.

On top of both, the inputs are FP16. A value drawn from an FP32 distribution is
rounded once on the way into the operand, before any multiply happens at all, and
that rounding is at the FP16 unit roundoff.

Putting the three together, the absolute error on one output element is bounded to
first order by

\begin{equation}
\left| \hat{c} - c \right| \lesssim
\left[ \frac{k}{16}\, u_{32} + u_{16} \right] \sum_{i} |a_i|\,|b_i|,
\label{eq:model}
\end{equation}

with $u_{32} = 2^{-24}$ and $u_{16} = 2^{-11}$.

The interesting part of Equation \ref{eq:model} is which term wins. At
$k = 4096$ the accumulation term is $256 \, u_{32} \approx 1.5 \times 10^{-5}$ and
the input rounding term is $u_{16} \approx 4.9 \times 10^{-4}$, so the input
rounding dominates by a factor of about 30. Deepening the accumulation, or
splitting K and reducing the partials, moves the smaller of the two terms. That is
why the accuracy of this kernel is essentially the accuracy of FP16 storage, and
why chasing accumulation order here would be effort spent on the term that is not
the problem.

\subsection{What is measured}

Three errors are recorded by \texttt{tests/test\_gemm\_tensor.cpp}, and each
answers a different question.

Against a same precision cuBLAS oracle the kernel agrees to about $10^{-7}$
relative Frobenius error. That is FP32 accumulate rounding and nothing else: the
kernel is numerically the same computation cuBLAS performs, not an approximation
of it. The correctness gate is $10^{-5}$, which the measured $10^{-7}$ clears by
two orders of magnitude, and every tile family member is held to the same gate at
unaligned shapes as well as aligned ones. The v1 gate was $5 \times 10^{-2}$,
which a kernel that silently dropped its last K stage would have passed; that is
why the number is what it is now.

Against an FP32 reference over the original data the measured error is about
$2.6 \times 10^{-4}$, which is the figure Equation \ref{eq:model} predicts once
the $u_{16}$ term is recognised as the dominant one. The suite records it as the
\texttt{vs\_fp32} property of the tensor test rather than as a number in a
document, so it is regenerated on every run and cannot drift away from the model
silently. The BF16 path measures about $2.1 \times 10^{-3}$, larger by roughly the
ratio the shorter mantissa predicts (8 explicit bits against FP16's 10).

\subsection{Which cuBLAS the oracle is}

The oracle only means something if the mode it ran under is stated, because with
FP16 inputs and \texttt{CUBLAS\_COMPUTE\_32F} cuBLAS is permitted to reduce
split-K partials in reduced precision, which moves both its speed and its accuracy
as a reference.

The tensor tests create a cuBLAS handle and do not call
\texttt{cublasSetMathMode}, so the oracle runs under the library default,
\texttt{CUBLAS\_DEFAULT\_MATH}, with reduced precision reduction permitted. That
is the honest statement of what the $10^{-7}$ agreement above was measured
against.

The benchmark path is stricter and newer. \texttt{bench\_all} sets the math mode
explicitly, reads it back with \texttt{cublasGetMathMode}, and records the
resulting name in the \texttt{cublas\_math\_mode} column of every row it writes,
so a row says for itself which oracle produced its baseline. The default is
still \texttt{CUBLAS\_DEFAULT\_MATH}; the
\texttt{-{}-disallow-reduced-precision-reduction} flag adds
\texttt{CUBLAS\_MATH\_DISALLOW\_REDUCED\_PRECISION\_REDUCTION} so the effect can be
measured rather than assumed.

The committed summary now answers this on its own face. All 982 measured rows carry
\texttt{cublas\_math\_mode = CUBLAS\_DEFAULT\_MATH}, so every percent of cuBLAS in
this report is against a vendor baseline that was permitted to reduce split-K
partials in reduced precision. That is the mode a caller gets without asking for
anything, which is why it is the one the headline is stated against, and it is the
more demanding denominator of the two: disallowing the reduced precision reduction
can only slow cuBLAS down, so quoting against the default is quoting against the
faster vendor. The stricter mode was not swept, so the size of that effect on this
part is not measured here.
